On dessine au tableau un polygone régulier à $n \geq 3$ côtés de côté $a > 1$ et on joue au jeu suivant.

Pierre dessine une diagonale (qui n'est pas un côté) du polygone, et Xavier efface un des deux demi-plans ouverts délimités par cette diagonale. On continue ainsi jusqu'à qu'il ne reste sur le tableau qu'un triangle dont on note les longueurs $x$, $y$ et $z$. On convient que Xavier gagne la partie si les nombres $x^2 - 1$, $y^2 - 1$ et $z^2 - 1$ sont encore les longueurs d'un triangle non aplati et qu'il la perd sinon.

Déterminer, en fonction de $n$, quel joueur possède une stratégie gagnante.
